\documentclass[11pt]{article}
\usepackage[T1]{fontenc}
\usepackage[utf8]{inputenc}
\usepackage{lmodern,microtype}
\usepackage[margin=0.88in]{geometry}
\usepackage{amsmath,amssymb,amsthm,mathtools}
\usepackage{booktabs,longtable,array,multicol}
\usepackage{xcolor,fancyhdr,enumitem}
\usepackage[colorlinks=true,linkcolor=blue!45!black,citecolor=blue!45!black,urlcolor=blue!45!black]{hyperref}
\hypersetup{
 pdftitle={Removing Factorial Authenticity: Binomial Prime-Power Certificates for Diophantine Consistency},
 pdfauthor={AI-assisted research prepared for Vladimir Reshetnikov and ProveIt},
 pdfsubject={An explicit quartic Diophantine equation equivalent to ZFC inconsistency},
 pdfkeywords={Diophantine equations, ZFC consistency, binomial coefficients, Chinese remainder theorem}
}
\usepackage{url}
\definecolor{ink}{RGB}{24,43,65}
\pagestyle{fancy}
\fancyhf{}
\fancyhead[L]{\small Factorial-free binomial CRT}
\fancyhead[R]{\small ProveIt research contribution}
\fancyfoot[C]{\thepage}
\setlength{\headheight}{14pt}
\setlength{\parskip}{0.35em}
\setlength{\parindent}{0pt}
\setlist{nosep,leftmargin=*}
\allowdisplaybreaks
\emergencystretch=2em
\newtheorem{theorem}{Theorem}[section]
\newtheorem{lemma}[theorem]{Lemma}
\newtheorem{proposition}[theorem]{Proposition}
\newtheorem{corollary}[theorem]{Corollary}
\theoremstyle{definition}
\newtheorem{definition}[theorem]{Definition}
\newtheorem{remark}[theorem]{Remark}
\newcommand{\N}{\mathbb N}
\newcommand{\Z}{\mathbb Z}
\newcommand{\Q}{\mathbb Q}
\newcommand{\Con}{\operatorname{Con}}
\newcommand{\lcm}{\operatorname{lcm}}
\newcommand{\vp}{v_p}
\newcommand{\PP}{\mathcal P}
\newcommand{\LL}{\mathcal L}
\input{results.tex}

\title{\color{ink}\bfseries Removing Factorial Authenticity\par
\large Binomial prime-power certificates for Diophantine consistency}
\author{AI-assisted mathematical research\smallskip\\
\normalsize Prepared for Vladimir Reshetnikov and the ProveIt project}
\date{25 October 2026}

\begin{document}
\maketitle
\thispagestyle{empty}
\begin{abstract}
We construct a fixed, fully emitted quartic polynomial over the unrestricted
integers whose unsolvability is provably equivalent in ZFC to
$\Con(\mathrm{ZFC})$, using the repository's numerical proof calculus
$C_{222}$. The principal endpoint has \NewVars{} variables,
\NewRows{} quadratic residuals, \NewTerms{} collected monomials, and
coefficient height \NewHeight{}. Every residual is printed in this report;
the accompanying JSON files also supply the complete collected polynomial.

The principal mathematical change is to remove the assertion that a
certificate variable equals $t!$. For a sufficiently large ordinary integer
$R$, the binomial coefficient $B=\binom{R+t}{t}$ itself retains a prime-power
selector for every one of $t$ rows. Its factors need not be pairwise coprime.
Conversely, every prime of $B$ can be assigned a compatible row, which closes
the completeness direction. A separate large-evaluation argument recovers
proof-label multiplicities without coprime numerical label weights. Thus an
explicit power replaces both uses of the old factorial. An odd-square radix,
a cofinal square choice for the beta base, and exact quadratic projections
complete the smaller equation.

The new arithmetic arguments are ordinary unbounded proofs. Finite programs
check the literal source edits, elimination identities, coordinate changes,
residual transformations, and every collected coefficient. The logical
$C_{222}$ calibration and unchanged Pell compiler are identified inherited
results. This release is not an end-to-end proof-assistant formalization,
a proof of consistency, or a claim of a global simplicity optimum.
\end{abstract}
\vfill
\begin{center}\small
Comparison snapshot: \texttt{3ecd6a6cf53d47519228328567f45fb7ee9781d9}.\\
All final unknowns in the principal equation range over $\Z$.
\end{center}
\newpage
\tableofcontents
\newpage

\section{The result and the meaning of ``simpler''}

The primary mathematical object is
\begin{equation}\label{eq:principal}
 \boxed{\quad P(x_0,\ldots,x_{\NewLast})
       =\sum_{j=0}^{\NewLastRow} R_j(x_0,\ldots,x_{\NewLast})^2=0,
       \qquad \boldsymbol x\in\Z^{\NewVars}.\quad}
\end{equation}
Appendix~\ref{sec:residuals} gives every coefficient of every quadratic
$R_j$. The delivered \nolinkurl{direct/principal/equation.json}
is the machine-readable definition; the neighboring
\nolinkurl{expanded_polynomial.json} is its completely collected expansion.
An intermediate register in the construction is a genuine variable only
when it remains in this final coordinate list.

\begin{theorem}[Consistency calibration]\label{thm:main}
For the fixed polynomial \eqref{eq:principal}, the new arithmetic proofs in
this report, composed with the pinned $C_{222}$/ZF/ZFC calibration, establish
\[
 \mathrm{ZFC}\vdash
 \left(\neg\exists\boldsymbol x\in\Z^{\NewVars}\;
 P(\boldsymbol x)=0\right)\ \longleftrightarrow\ \Con(\mathrm{ZFC}).
\]
No consistency assumption is used to prove this equivalence.
\end{theorem}

The direct route compiles numerical proof rules. It does not simulate
K215, specialize Jones's universal polynomial, or contain a large hidden
machine coefficient. The coefficients are the ordinary small integers
printed in the residual appendix. This distinction matters for inspection
and formalization, even though existing native constructions use many fewer
variables.

\subsection{Exact comparison}

\input{comparison.tex}

The baseline is the latest common-remainder direct quartic in the eight
incoming submissions present at the comparison commit~\cite{common,repo}.
The separate centered-residue and residual-basis submissions supply useful
independent transformations~\cite{centered,bases}. We credit these
transformations rather than presenting their composition as a new theorem.
The factorial-free source is a change to the arithmetic architecture, and
is the principal contribution of this report.

There is no metric-independent smallest equation. The inspected repository
already has nine-positive-variable equations, ten-integer equations using a
different reduction, and a fifteen-positive-variable 66-operation circuit.
Those are different tradeoffs in domain, degree, coefficient description,
and evaluation cost. Our improvement concerns the direct small-coefficient
quartic family. The lower-dimensional dense companion deliberately spends
more support to remove one coordinate. Neither a global minimum nor
literature-wide priority is asserted.

\section{Why the old factorial can be removed}

In the inherited direct architecture one integer $F=t!$ serves two jobs.
First, the weights $1+jF$, for labels $1\le j\le t$, are pairwise coprime,
so their prime valuations recover proof-label multiplicities. Second,
$\binom{F^2+t}{t}$ decomposes into pairwise coprime reciprocal factors
$1+F^2/i$, providing a CRT slot for each proof row. Certifying $F=t!$
requires a self-referential ceiling identity and its supporting bounds.

We replace the two jobs independently. The label product contains only
boundedly many factors, so equality at a sufficiently large integer can be
lifted to an identity of polynomials. For row selection, pairwise coprime
reciprocal factors are stronger than needed: appropriate prime powers
inside one binomial coefficient already provide all row selectors. The
remaining issue is completeness at \emph{every} prime of that coefficient;
Lemma~\ref{lem:all-primes} handles it explicitly.

The centered-residue predecessor already considered coefficient separation
for the label product, but identified the independent factorial dependence
of the CRT carrier as an obstacle~\cite{centered}. The prime-power compiler
below removes that remaining dependence. This is the substantive new
architectural step.

The new arithmetic choice is
\begin{equation}\label{eq:power-F}
 d=2(t+1),\qquad F=d^t=2^t(t+1)^t.
\end{equation}
The product argument below first recovers equal factor counts, so the
smaller bound $F>(t+1)^t$ suffices.
This is a semantic description of an ordinary polynomial certificate.
The old power request at index $t$ is reused, so the number of power
requests remains eighteen. Variable exponentiation is not an operation
in the final polynomial.

\section{Prime powers in a binomial modulus}

Set $\Lambda_t=\lcm(1,\ldots,t)$ for $t\ge1$. Write $v_p(n)$ for
the exponent of the prime $p$ in a positive integer $n$.

\begin{lemma}[A classical denominator bound]\label{lem:lcm}
For $1\le i\le t$,
\[
 t\binom{t-1}{i-1}\mid\Lambda_t.
\]
\end{lemma}
\begin{proof}
For $d=p^a$, the summand in the factorial valuation of the left side is
\[
 \left\lfloor\frac{t}{d}\right\rfloor
 -\left\lfloor\frac{i-1}{d}\right\rfloor
 -\left\lfloor\frac{t-i}{d}\right\rfloor.
\]
It is zero or one, because $(i-1)+(t-i)=t-1$, and it vanishes for $d>t$.
The valuation is therefore at most $\lfloor\log_p t\rfloor$, exactly
the valuation of $\Lambda_t$.
\end{proof}

This divisibility is classical and follows also from the lcm identity for
binomial rows discussed by Farhi~\cite{farhi}. Its use below is credited
accordingly; it is not claimed as a new number-theoretic identity.

\begin{lemma}[Every prime has a compatible row]\label{lem:all-primes}
Let $R,t\ge1$, $B=\binom{R+t}{t}$, and $p^e\Vert B$. There is an
$i\in\{1,\ldots,t\}$ for which $p^e\mid R+i$.
\end{lemma}
\begin{proof}
Let $v=\max_{1\le i\le t}v_p(R+i)$. Legendre's formula gives
\[
 e=\sum_{a\ge1}\left(
 \left\lfloor\frac{R+t}{p^a}\right\rfloor
 -\left\lfloor\frac{R}{p^a}\right\rfloor
 -\left\lfloor\frac{t}{p^a}\right\rfloor\right).
\]
Each summand is zero or one. Above $a=v$ the interval contains no
multiple of $p^a$; necessarily $p^a>t$, since any $t$ consecutive
integers contain such a multiple if $p^a\le t$. All those summands
are zero, so $e\le v$. A row attaining the maximum works.
\end{proof}

\begin{lemma}[Every row has a sufficiently large selector]\label{lem:each-row}
Let $R\ge\Lambda_t^2$ and $B=\binom{R+t}{t}$. For each
$i\in\{1,\ldots,t\}$ there is a prime power $p^e\Vert B$ such that
\[
                 p^e\mid R+i,\qquad p^e>t.
\]
\end{lemma}
\begin{proof}
Because $R+i>\Lambda_t^2$, the integer $R+i$ does not divide
$\Lambda_t^2$. Hence some prime $p$ has
$v=v_p(R+i)>2a$, where $a=v_p(\Lambda_t)=\lfloor\log_p t\rfloor$.
In particular $p^v>t$. For $j\ne i$, the valuation of $R+j$ is
therefore the valuation of $j-i$. Put
\[
 d=v_p\!\left(t\binom{t-1}{i-1}\right).
\]
The valuation of $B$ is exactly
\[
 e=v+v_p((i-1)!(t-i)!)-v_p(t!)=v-d.
\]
Lemma~\ref{lem:lcm} gives $0\le d\le a$. Consequently
$e\ge v-a>a$, so $p^e>t$, while $e\le v$ gives $p^e\mid R+i$.
This also covers $p>t$, for which $a=d=0$.
\end{proof}

The two lemmas have different purposes. A row selector proves soundness
for that row. Lemma~\ref{lem:all-primes} supplies a compatible assignment
at every prime of $B$ in the converse construction, including primes
that were not chosen as row selectors. Omitting the latter step would
leave a real completeness gap.

\subsection{A small example}

Take $t=3$, $R=37$. Here $\Lambda_3^2=36$ and
\[
 B=\binom{40}{3}=9880=2^3\cdot5\cdot13\cdot19.
\]
The consecutive factors are $38,39,40$, which are not pairwise coprime.
Their full prime powers in $B$ can be assigned as follows:
\begin{center}
\begin{tabular}{ccc}
\toprule
Prime power in $B$ & Compatible factor & Row index $i$\\
\midrule
$19$ & $38$ & $1$\\
$13$ & $39$ & $2$\\
$8$ and $5$ & $40$ & $3$\\
\bottomrule
\end{tabular}
\end{center}
Every row has a selector greater than $3$, and all primes of $B$ have
a compatible row. No factorization into $1+F^2/i$ is used.

\section{A bounded-row certificate without factorial authenticity}

We first recall the exact one-sided range fact used in the repository.
\begin{lemma}[Binomial range]\label{lem:range}
If $p$ is prime, $A\in\N$, $k,e\ge1$, and $p^e\mid\binom{A}{k}$, then
$A\bmod p^e<k$. The convention is $\binom{A}{k}=0$ for $A<k$.
\end{lemma}
\begin{proof}
The case $A<k$ is immediate. Otherwise, if none of the $k$ consecutive
integers $A-k+1,\ldots,A$ is divisible by $p^e$, none is divisible by
a higher power either. The valuation formula for the binomial is a sum
of zero-one terms, with all terms from level $e$ onward zero. The total
is at most $e-1$, a contradiction. Thus $A-r$ is divisible by $p^e$
for some $0\le r<k$; its ordinary residue is at most $r$ and hence
is less than $k$.
\end{proof}

\begin{theorem}[Binomial CRT compiler]\label{thm:crt}
Let $t\ge1$, $R\ge\Lambda_t^2$, $B=\binom{R+t}{t}$, and $N=B^K$.
Let $b_1,\ldots,b_d\in\N$, with
$K\ge\max(t,b_1+1,\ldots,b_d+1,1)$.
Suppose a fixed integer polynomial $\LL$ satisfies
\[
 0\le\LL(i,\boldsymbol w)<2^K
 \quad(0\le i<t,\ 0\le w_j\le b_j).
\]
Then the following are equivalent:
\begin{enumerate}
\item Every row $i<t$ has a bounded tuple $\boldsymbol w$ with
      $\LL(i,\boldsymbol w)=0$.
\item There are nonnegative $C,A_1,\ldots,A_d$ satisfying
\begin{align}
 B&\mid C+R+1,\label{eq:raw-index}\\
 N&\mid\binom{C}{t},\label{eq:range-index}\\
 N&\mid\binom{A_j}{b_j+1}\quad(1\le j\le d),\label{eq:range-columns}\\
 N&\mid\LL(C,\boldsymbol A).\label{eq:local-congruence}
\end{align}
\end{enumerate}
In the forward construction the representatives may all be chosen in
$[1,N^2]$. Soundness does not require that upper bound.
\end{theorem}

\begin{proof}
For soundness fix $i\in\{1,\ldots,t\}$ and choose $p^e$ from
Lemma~\ref{lem:each-row}. By Lemma~\ref{lem:range}, the index range
gives $C\equiv s\pmod{p^{Ke}}$ with $0\le s<t$. Reduction of
\eqref{eq:raw-index} modulo $p^e$ gives
$s\equiv i-1\pmod{p^e}$. Since $p^e>t$, we have $s=i-1$ exactly.
Each other range condition supplies $A_j\equiv w_j\pmod{p^{Ke}}$
with $0\le w_j\le b_j$. Therefore $p^{Ke}$ divides
$\LL(i-1,\boldsymbol w)$. The latter lies in $[0,2^K)$ and
$p^{Ke}\ge2^K$, so it is zero.

Conversely choose a valid tuple for each row. For \emph{every} prime
$p\mid B$, write $e=v_p(B)$ and choose a compatible $i$ by
Lemma~\ref{lem:all-primes}. Prescribe, modulo $p^{2Ke}$,
\[
 C\equiv i-1,\qquad A_j\equiv w_j^{(i-1)}.
\]
These moduli are pairwise coprime and have product $N^2$. CRT supplies
positive representatives at most $N^2$. The raw index and local
congruences follow. For a range lower index $k\le K$, its selected
residue $w$ is in $[0,k-1]$. If $A<k$ its binomial is zero; otherwise
the falling factorial for $\binom{A}{k}$ has the factor $A-w$ divisible
by $p^{2Ke}$. Division by $k!$ loses at most
$v_p(k!)\le k-1<K\le Ke$. Thus the binomial is divisible by
$p^{Ke}$. This holds for every prime of $B$ and for every range.
\end{proof}

\paragraph{Why a size hypothesis cannot simply be dropped.}
Let $t=2$, $R=1$, $B=3$, $K=2$, $N=9$, $C=1$, and
$\LL(i)=(i-1)^2$. Then all three applicable global congruences hold:
$3\mid C+2$, $9\mid\binom C2=0$, and $9\mid\LL(C)=0$.
But row $0$ is false. The missing row's prime factor was cancelled by
the binomial denominator. This explicit counterexample is included in
the finite regression script.

\section{Proof-label balance from a polynomial identity}

Let a calibrated Horn rule have one conclusion and at most two premises.
Pad unused premise slots by the freely derivable truth code $0$. Suppose
there are $t$ rows and all their labels lie in $[0,t]$. For a fixed
positive integer $F$ put
\[
 W_j=1+jF,\qquad W_0=1.
\]

\begin{lemma}[Large evaluation recovers multiplicities]\label{lem:labels}
If $F>(t+1)^t$, then
\begin{equation}\label{eq:label-product}
 \prod_{i<t}W_{v_i}=(1+F)\prod_{i<t}W_{p_i}W_{q_i}
\end{equation}
implies, at each positive label $j$, that its conclusion multiplicity
minus its premise multiplicity is $1$ for $j=1$ and $0$ otherwise.
\end{lemma}
\begin{proof}
First discard all unit factors with label zero. The left side has at
most $t$ nonunit factors; the right side has at most $2t+1$. If these
counts differ, their minimum $r$ is at most $t$. The side with more
factors is at least $(F+1)^{r+1}$, whereas the side with fewer is at
most $(1+tF)^r\le t^r(F+1)^r$. Since $F+1>t^t\ge t^r$, this is
impossible. Thus the factor counts agree and both are at most $t$.

Replace $F$ by an indeterminate $z$. Each side now has nonnegative
coefficients and coefficient sum at most $(t+1)^t<F$. Thus every
coefficient of the difference has absolute value less than $F$.

A nonzero integer polynomial with that coefficient bound cannot vanish
at $F$: take its lowest nonzero coefficient $a_r$, divide its evaluation
by $F^r$, and reduce modulo $F$. This would force $F\mid a_r$ with
$0<|a_r|<F$. Therefore \eqref{eq:label-product} is an identity in
$\Z[z]$. Distinct positive labels give distinct irreducible linear
factors $1+jz$ over $\Q$, with roots $-1/j$. Unique factorization
gives the required multiplicities. Label $0$ contributes only the unit.
\end{proof}

Pairwise coprimality of the \emph{evaluated} integers $W_j$ is not
required. The choice \eqref{eq:power-F} has the needed strict bound.

\begin{lemma}[Single-conclusion flow]
If the balance in Lemma~\ref{lem:labels} is supplied by valid
single-conclusion Horn rules and truth is derivable, then label $1$ is
derivable.
\end{lemma}
\begin{proof}
Take the finite closure of truth under the finitely many rule instances
in the certificate. If $1$ is outside this closure, every rule whose
conclusion is outside has at least one premise outside. Summing
conclusions minus premises over outside labels is therefore at most
zero. The balance makes that sum one. This contradiction gives a finite
derivation. Repeated premises and cyclic rule multisets cause no problem.
\end{proof}

This last lemma is the inherited commutative-certificate argument, included
to make clear why an unordered balance does not permit circular false
proofs~\cite{flow}. The single-conclusion restriction is essential.

\section{The selected arithmetic front}

All expressions in this section describe values of explicitly certified
polynomial registers. The final system itself is entirely quadratic.
Let $t=16+h^2$ and retain the same $C_{222}$ local rule polynomial, range
columns, semantic anchor, and eighteen power requests as the parent.
The selected changes are
\begin{align}
 d&=2(t+1),&F&=d^t,\label{eq:new-F}\\
 s&=(2v+1)^2,&y&=Fs,\label{eq:new-y}\\
 q&=(2z+1)^2,&H&=2+C+\sum_{j=1}^{8}A_j.\label{eq:new-q}
\end{align}
Here $C$ and the eight $A_j$ are the actual CRT representatives; some
stored source coordinates are these values minus one. The distinction
is retained explicitly in the mutation ledger. The former factorial
ceiling rows and $y=V_0(V_0-1)$ are deleted. The variable $x$ is unrestricted;
there is \emph{no} parity equation for $x$ in the selected source.

The unchanged endpoint and arithmetic requests include
\begin{align}
 x&=F+1+(y+1)u_0,\label{eq:first-end}\\
 x&=1+[1+(t+1)y]u_t,\label{eq:last-end}\\
 B h_B&=C+F^2+1,\label{eq:index-row}\\
 N&=B^x,&q&=(BH)^x.\label{eq:two-powers}
\end{align}
The digit called the ``factorial digit'' in old filenames now has only
the intended meaning
\begin{equation}\label{eq:new-binomial}
                     B=\binom{F^2+t}{t}.
\end{equation}
It no longer authenticates a factorial. Old variable names retained in
data provenance do not change this mathematical interpretation.

\subsection{The unchanged Pell interface and signed bases}

For an integer center $A$, let $X_k(A),Y_k(A)$ be the Pell polynomials
defined by
\[
 X_0=1,\quad Y_0=0,\qquad
 X_{k+1}=AX_k+(A^2-1)Y_k,\quad Y_{k+1}=X_k+AY_k,
\]
and $X_{-k}=X_k$, $Y_{-k}=-Y_k$. The retained signed square frame,
semantic anchor, and pooled small-residue proof establish these exact
coordinates at every requested signed index. They bound every actual
base, index, and output atom in absolute value strictly below $W$, with
$W\ge258$ and $|A|>W^W$. These are properties of the unchanged graph;
the new base $2(t+1)$ is explicitly inserted into its atom list.
The relevant ordinary proofs are retained in \nolinkurl{direct/parent/}
and identified in the source ledger~\cite{common,centered}.

For a nonzero signed base $n$ define $T_n=2An-n^2-1$. The certified
power congruence and the recurrence imply
\[
 m\equiv n^k\pmod{T_n}\quad(k\ge0),\qquad
 mn^{-k}\equiv1\pmod{T_n}\quad(k<0).
\]
The modulus satisfies $|T_n|>W^W+W$. For $|n|\ge2$ the corresponding
absolute errors are less than the modulus. A negative $k$ would force
the impossible integer equality $mn^{-k}=1$. Thus $k\ge0$ and
$m=n^k$. For $n=\pm1$ the same recurrence gives exactly
$m=(\pm1)^k$ for signed $k$. A base of zero is not covered by this
argument, and is excluded explicitly below when needed.

\section{Soundness without a circular sign argument}\label{sec:soundness}

The order of recovery is essential. In particular, neither the sign of
$B$ nor the meaning of its binomial digit is assumed in an earlier step.

\subsection{Recover the power, representatives, and positive exponent}

The request \eqref{eq:new-F} has known positive base and $t\ge16$, so
it gives
$F=(2(t+1))^t>2^{32}$ independently of every digit. Equation
\eqref{eq:new-y} gives $y\ge F>0$. Equation \eqref{eq:new-q} gives
$q\ge1$, so every numerator request at base $q+1\ge2$ forces its
actual upper index to be nonnegative. Hence $C,A_j\ge0$ and $H\ge2$.
The raw index row \eqref{eq:index-row} excludes $B=0$ because its
right side is positive. Therefore $|BH|\ge2$.

The signed-base request $q=(BH)^x$ now forces $x\ge0$ and the actual
power equality, even if the tentative value of $B$ is negative. The case
$x=0$ is impossible in \eqref{eq:last-end}: its coefficient
$1+(t+1)y$ is greater than one and cannot divide $-1$. Thus $x>0$.
If $u_0\le-1$, \eqref{eq:first-end} gives
$x\le F-y\le0$, also impossible. Consequently
\begin{equation}\label{eq:x-bound}
             x\ge F+1\ge3,
       \qquad x\ge2+(t+1)y\ge2+(t+1)F.
\end{equation}
The second inequality follows from $u_t\ge1$ in the last endpoint.

Since $q>0$ and $q=(BH)^x$, we now know
\begin{equation}\label{eq:q-growth}
 q=|BH|^x\ge2^x>1,\qquad |B|\le |BH|<q/2.
\end{equation}
All powers at base $q$ are therefore genuine, and $t,t+1,y+1\le x$.
The exact digit for $B$ uses no range factor $N$, so the latter's positive
interpretation can wait until $B>0$ has been proved.

\subsection{Authenticate the binomial from small residues}

For $0\le j\le t$ a convenient elementary bound is
\[
 \binom{F^2+t}{j}\le(F^2+t)^t\le(F^2+F)^t
                         \le2^{t(F-1)}.
\]
Here $F\ge t$ and $F^2+F\le2^{F-1}$ for $F\ge7$, by a direct
base case and induction. Combining with \eqref{eq:x-bound} gives
\[
 \frac{2^{t(F-1)}}{q}\le2^{-F-t-2},\qquad
 \frac{2^{t(F-1)}}{q-1}<\frac1{32}.
\]
Thus the actual low remainder of $(q+1)^{F^2+t}$ modulo $q^t$ is
less than $q^t/32$, and the coefficient
$C_* = \binom{F^2+t}{t}$ is less than $q/32$.

The retained common-remainder construction uses $M=q^x$. For each
digit denominator $D_i\mid M$, it replaces a remainder $r_i$ by
$\rho_i=(M/D_i)r_i$ and imposes
\begin{equation}\label{eq:common-norm}
 M^2=1+4\sum_{i=1}^{11}\rho_i^2+\alpha^2+\beta^2+\gamma(\gamma+1).
\end{equation}
The last three terms are nonnegative for integers, so $|r_i|<D_i/2$.
The scaled digit row makes $r_i$ an integer congruent to its actual
remainder. If that actual remainder is in $[0,D_i/32)$, the difference
has absolute value below $D_i$, and therefore is zero. This also explains
why the normalizing ratio variables do not introduce rational solutions.

Apply this first to the exact digit for $B$. It gives
$B\equiv C_*\pmod q$. By \eqref{eq:q-growth},
$|B-C_*|<q/2+q/32<q$. Therefore
\[
                         B=C_*>0.
\]
Only at this point do we use $B\ge2$, $N=B^x>0$, and the usual range
interpretation. This signed reconstruction replaces the old output-slack
positivity guard.

\subsection{Authenticate every remaining range}

For an ordinary upper representative $a\ge0$ and a required index
$0\le j\le k\le x$,
\[
 \binom aj\le\max(1,a)^k\le H^x.
\]
This includes $a=j=0$, a case that an $a^k$ bound would miss. The
derived label has the inherited bound $h_*\le H^{\ell}$ with
$\ell=4194304$. The constant shift in $H$ still gives the required
bound for zero representatives. The growth in \eqref{eq:x-bound}
implies $\ell(t+1)<x$, so the derived-label coefficients also lie below
$H^x$. Because $q=B^xH^x$ and $B\ge2$,
\[
                  \frac{H^x}{q-1}<\frac2{2^x}<\frac1{32}.
\]
Every remaining low remainder is therefore selected correctly by
\eqref{eq:common-norm}. The independent-column divisibilities are those
of Theorem~\ref{thm:crt}. There is also a range for a polynomially derived
label $J(\boldsymbol A)$, rather than a new independent column.
For a selected prime power, that range gives
$J(\boldsymbol A)\equiv r\pmod{p^{xe}}$ with $0\le r\le t$.
The column congruences imply
$J(\boldsymbol w)\equiv r\pmod{p^{xe}}$.
The inherited local ledger bounds this nonnegative derived value strictly
below $2^x\le p^{xe}$; hence $J(\boldsymbol w)=r\le t$.
This is the precise extension of the compiler used for the derived label.

The unchanged local majorant is
$0\le\LL<\mathcal R^E$, where
$\mathcal R=2+t+x+y+F$ and $E=8388608=2^{23}$.
From $y\ge F$ and \eqref{eq:x-bound}, $\mathcal R<2x$.
Also $x>2^{32}$, and the elementary dyadic estimate
$2^{23}\log_2(2x)<x$ gives
\[
                  \LL<(2x)^E<2^x.
\]
The exponent $K=x$ therefore satisfies every range and evaluation budget.
Since $F>t!\ge\Lambda_t$, the radius $R=F^2$ satisfies the new
CRT hypothesis. Every local row is recovered.

\subsection{The boundary case \texorpdfstring{$y=F$}{y = F}}

The cofinal square parametrization permits $y=F$ algebraically. It must
not be silently treated as the old strict bound. After bounded row
decoding, the first endpoint has remainder zero modulo $y+1$, because
$F+1=y+1$. Thus the initial accumulator is zero. Each local update is
\[
 S_{i+1}W_{v_i}=S_iW_{p_i}W_{q_i},
\]
with every weight positive, so zero propagates to the final accumulator.
The last endpoint has ordinary remainder one and forces that final
accumulator to be one, a contradiction. Hence every actual zero has
$y>F$. Now $F+1\le y$, and the first endpoint restores the intended
initial accumulator $S_0=F+1$.

All accumulators are positive by induction. Multiplying their updates
and cancelling gives \eqref{eq:label-product}. Lemma~\ref{lem:labels}
and the single-conclusion flow argument yield a finite $C_{222}$
refutation. The retained logical calibration then yields a finite
syntactic ZFC refutation.

\section{Constructive completeness}

Start with a finite calibrated $C_{222}$ refutation. Expand it to a finite
proof tree, and choose a horizon $t=16+h^2$ beyond its size, all its
schema parameters, and all required local label sums. This choice precedes
all large arithmetic witnesses. Define $F$ by \eqref{eq:power-F}, and
construct the pending-multiset accumulators using weights $1+jF$.
Truth padding gives exactly $t$ rows.

Let $O$ be the odd part of $t!$. Choose an arbitrarily large odd multiple
$w$ of $O$ and set $y=Fw^2$. Because
\[
 v_2(F)=t\bigl(1+v_2(t+1)\bigr)\ge t>v_2(t!),
\]
$y$ is an even multiple of $t!$. Choose $w$ large enough that $y>F$
and every accumulator is at most $y$. The beta moduli
$1+(i+1)y$, $0\le i\le t$, are then pairwise coprime: a common
prime would divide $j-i\le t$ and hence $y$, while also dividing
$1+(i+1)y$, impossible. They are all odd. CRT supplies a common beta
code $x$; adding their product changes its parity, so choose $x$ even.
The genuine endpoint quotients satisfy \eqref{eq:x-bound}, and the local
beta quotients are strictly below $x$ because their moduli are at least two.

Define $R=F^2$, $B=\binom{R+t}{t}$, and $N=B^x$. The new binomial
CRT construction supplies all actual representatives modulo $N^2$.
The integer $B$ is odd: for $1\le i\le t$,
$v_2(R)>\lfloor\log_2t\rfloor$ gives $v_2(R+i)=v_2(i)$, so their
valuations cancel exactly against $t!$.

Consequently the representative classes modulo $N^2$ have free parity.
Choose the dominant schema representative $a$ odd, and the other eight
representatives even. Take $a$ in $[2N^2,4N^2)$ by adding suitable
multiples of the odd modulus. The inherited shell constructors give
$h_*=1+a+b\pmod2$ and $h_*\ge2a^2$, so the semantic anchor is even
and dominates the other indices. More explicitly, choose the other eight
positive representatives at most $2N^2$. The front indices satisfy
$F^2+t\le B<N^2$ and $x,y+1,t<N^2$, whereas
$h_*\ge8N^4>2\max(\text{other indices})$.
The derived-label range is preserved because polynomial evaluation
respects the CRT congruences modulo $p^{2xe}$, and its denominator loses
at most $v_p((t+1)!)\le t<x$ powers of $p$.
Meanwhile $H=2+\sum a_j$ is odd.
Since $B$ is odd and $x$ is even,
\[
                  q=(BH)^x
\]
is an odd square and has the parametrization in \eqref{eq:new-q}.
All four denominator groups divide $q^x$. The genuine normalized
remainders satisfy $0\le\rho_i<q^x/32$, so
\[
 4\sum_{i=1}^{11}\rho_i^2<\frac{44}{1024}(q^x)^2<(q^x)^2.
\]
The nonnegative integer gap in \eqref{eq:common-norm} is represented by
$\alpha^2+\beta^2+\gamma(\gamma+1)$, using the standard three-square
representation retained in the parent. Explicitly, for any gap $g\ge0$,
Legendre's three-square theorem represents $4g+1$; reduction modulo four
makes exactly one of its three squares odd. Writing the even roots as
$2\alpha,2\beta$ and the odd root as $2\gamma+1$ gives the stated form.

The unchanged constructive Pell frame and anchor supply the finite
family of genuine powers. The new request has index $t$, which was
already present; no new index-compatibility premise has appeared.
This constructs a zero of the changed arithmetic source. Exact integral
projections and reversible row/coordinate operations transport it to
\eqref{eq:principal}. Completeness uses factorials in the \emph{proof of
existence of a suitable multiple}; no unknown is required to equal a
factorial in the certificate.

\section{From the arithmetic source to a literal quartic}

The new source keeps the calibrated local $C_{222}$ polynomial, its
eighteen power slots, and the common-remainder transformation. All new
definitions are quadratic register rows. The single variable power in
\eqref{eq:new-F} is compiled by the retained Pell graph at the existing
index $t$. Thus every final residual is an ordinary integer polynomial
of degree at most two; summing their squares gives degree at most four.
The emitted expansion has nonzero degree-four coefficients, so its
degree is exactly four.

The compiler performs three kinds of finite transformations. First, it
substitutes variables whose defining row is monic and whose substitution
keeps every remaining row quadratic. The deleted variable is restored
uniquely by its integer polynomial. Second, it exposes three already
valid affine consequences by exact row identities and projects their
variables. A fourth such projection gives the denser companion. Third,
it applies recorded integral coordinate translations and elementary
residual-row operations. A row operation replacing $R_i$ by
$R_i+nR_j$ is reversed by subtracting $nR_j$; it preserves the common
zero set even though it changes the sum-of-squares polynomial.

These facts are checked coefficient by coefficient. The row-reduction
search is a heuristic used only to find a useful basis; its frozen finite
operation list is the certificate. Reproducing that list does not require
rerunning the search. No random modular agreement is treated as an
identity proof.

\section{Logical calibration and verification boundary}

The logical target is the repository's fixed numerical calculus
$C_{222}$ with truth code $0$ and falsity code $1$. The retained named
interfaces are
\begin{center}\small
\nolinkurl{KEndpoint.c222_iff_zf_refutation},\\
\nolinkurl{ConstructibleChoice.zf_refutation_iff_zfc_refutation}.
\end{center}
The first calibrates the numerical rules with syntactic ZF refutability;
the second uses constructible choice for ZF/ZFC equiconsistency. Ordinary
proof sources and the relevant interface source files are retained with
their provenance. Their complete transitive proof-assistant closure was
not rebuilt in this task.

The new proofs above establish, in finite arithmetic formalizable in ZFC,
both implications between a zero of the emitted polynomial and a finite
$C_{222}$ refutation. Composing the identified calibration gives
\[
 \mathrm{ZFC}\vdash
 \bigl(\exists\boldsymbol x\;P(\boldsymbol x)=0\bigr)
 \longleftrightarrow
 \bigl(\exists n\;\operatorname{Proof}_{\mathrm{ZFC}}(n,\bot)\bigr).
\]
Negation gives Theorem~\ref{thm:main}. This is an ordinary mathematical
provability argument, not an emitted object-language proof string. The
distinction between an unbounded proof and a finite artifact check is
preserved throughout the release.

\input{verification.tex}

\section{What is new, inherited, and still open}

The new contribution is the explicit replacement of factorial
authenticity by a binomial prime-power compiler with two-sided row coverage,
together with the compatible proof-label and signed-source changes.
The sharpened $\Lambda_t^2$ threshold is proved here using a classical
lcm divisibility. The polynomial-evaluation device itself is elementary;
its use with the bounded $C_{222}$ label product removes the need for
coprime evaluated weights. The cofinal odd-square beta base adds a further
semantic saving, with the $y=F$ exceptional branch explicitly excluded.

The numerical calculus, signed Pell frame, range lemma, common-remainder
normalization, and affine residual exposures have identified predecessors.
The archive records their exact source boundaries. This report does not
claim that an exhaustive literature search establishes worldwide novelty.
It establishes a complete new arithmetic source relative to the inspected
repository and an exact smaller direct quartic.

Several questions remain worthwhile. The selector threshold
$R\ge\Lambda_t^2$ is sufficient and need not be optimal. A row-specific
threshold using $\Lambda_t\,t\binom{t-1}{i-1}$ may reduce it further,
although this would not presently lower the emitted graph size. The
largest remaining arithmetic cost is the shared Pell certification, so a
new finite equation for its simultaneous powers would have greater impact
than another local row-basis search. Finally, an end-to-end formalization
should start with the prime-power coverage theorem and signed bootstrap:
they are compact statements with explicit exceptional cases and clear
interfaces to the existing graph.

\section{Files, reproduction, and integration}

The archive is intended as a new contribution under
\nolinkurl{Computability/DiophantineConsistency/con-zfc-diophantine-forms},
or initially in \nolinkurl{docs/incoming}. It includes this source and its
compiled PDF, the literal principal and companion equations, all collected
coefficients, construction and checking programs, frozen transformation
certificates, mathematical reviews, retained ordinary proof dependencies,
licenses, and a checksum manifest.

Run \texttt{python3 -B verify.py} from the archive root for the selected
finite checks. They use Python's standard library and do not need a network
or a consistency oracle. The evaluator accepts a supplied integer tuple;
it does not attempt a root search. Rebuild this report with two runs of
\texttt{pdflatex -interaction=nonstopmode -halt-on-error main.tex}
from the \nolinkurl{report/} directory. Different TeX versions can change
PDF bytes without changing the polynomial or mathematics.

\begingroup
\raggedright
\begin{thebibliography}{9}
\bibitem{repo} V. Reshetnikov and ProveIt contributors,
\emph{ProveIt: DiophantineConsistency}, comparison commit
\texttt{3ecd6a6cf53d47519228328567f45fb7ee9781d9}, 25 October 2026.
\url{https://github.com/VladimirReshetnikov/ProveIt/tree/3ecd6a6cf53d47519228328567f45fb7ee9781d9/Computability/DiophantineConsistency}.

\bibitem{common} ProveIt incoming research,
\emph{Common remainder certificates for Diophantine consistency},
\nolinkurl{ConZFC_Common_Remainders_Radix47_2026-10-25.zip},
\nolinkurl{direct/proofs/NORMALIZED_REMAINDERS_PROOF.md} and its retained
orientation/strict-budget proof sources. Available in \texttt{docs/incoming}
at the comparison commit; exact hashes are in this archive's source ledger.

\bibitem{centered} ProveIt incoming research,
\emph{Centered residue certificates and reusable Pell denominators},
\nolinkurl{ConZFC_Centered_Residues_2026-10-25.zip}, 25 October 2026.
The odd-square radix and compatible parity construction are predecessors
to the changed bootstrap in this report.

\bibitem{bases} ProveIt incoming research,
\emph{Residual bases and denominator engineering},
\nolinkurl{ConZFC_Residual_Bases_Denominator_Engineering_2026-10-25.zip},
25 October 2026. Source of the exact affine residual exposures used here.

\bibitem{flow} ProveIt direct-calculus source,
\emph{Commutative refutation certificates}; retained
\nolinkurl{COMMUTATIVE_CERTIFICATES.md}, numerical-calculus proofs and
calibration interfaces in the source closure of \cite{common}.

\bibitem{farhi} B. Farhi, \emph{An identity involving the least common
multiple of binomial coefficients and its application}, 2009,
\href{https://arxiv.org/abs/0906.2295}{arXiv:0906.2295v2}.

\bibitem{universal} J. Bayer, M. David, M. Hassler, Y. Matiyasevich,
and D. Schleicher, \emph{Diophantine equations over $\Z$: universal bounds
and parallel formalization},
\href{https://arxiv.org/abs/2506.20909}{arXiv:2506.20909v3},
22 October 2026. Background on the domain/degree/variable tradeoff and the
distinction between ordinary and formal verification.
\end{thebibliography}
\endgroup

\appendix
\section{Complete quadratic residuals}\label{sec:residuals}

All coordinates below are unrestricted integers. An empty monomial in the
JSON definition denotes the constant one; repeated coordinate indices
denote repeated factors. The list below is generated from the final
canonical JSON, and the release verifier checks their agreement. The
single equation is exactly the sum of the squares of these residuals.

\input{residuals.tex}

\end{document}

